\section{Comparison and Evidence Plan}
\label{sec:comparison}

\subsection{Current Evidence Status}

The implementation has unit coverage for parsing, deterministic behavior,
seed-prefix separation, direction selection on simple grids, empty-centroid
failure, and full \texttt{run\_all\_splits\_mka} smoke behavior. That is enough
to trust the algorithmic path for development. It is not enough to publish
statewide comparative performance claims.

\begin{table}[t]
\centering
\small
\caption{Evidence status for Moving-Knife claims.}
\label{tab:evidence-status}
\renewcommand{\arraystretch}{1.25}
\begin{tabular}{@{}p{0.28\linewidth}p{0.20\linewidth}p{0.40\linewidth}@{}}
\toprule
Claim & Status & Required next artifact \\
\midrule
CLI parses moving-knife & implemented & keep regression tests \\
Centroids are required & implemented & manifest error examples \\
Orientation sweep is deterministic & implemented & per-node angle ledger \\
Centroid-MEC Reock ranking works & implemented & replayable trace package \\
Polygon-boundary Reock & not implemented here & boundary-vertex scorer \\
Polsby-Popper MKA & enum only & perimeter-aware split scorer \\
NC/FL/WA superiority & target claim & archived same-metric runs \\
Partisan neutrality & target claim & multi-seed outcome packages \\
\bottomrule
\end{tabular}
\end{table}

\subsection{How the Benchmark Should Look}

The old version of this paper reported concrete NC/FL/WA numbers. Those
numbers should be treated as placeholders until the run directories, inputs,
seeds, manifests, and score scripts are archived. The correct benchmark table
has this shape:

\begin{table}[h]
\centering
\caption{Moving-Knife benchmark package target. All entries are pending
archived same-metric runs.}
\label{tab:reock-comparison}
\renewcommand{\arraystretch}{1.25}
\begin{tabular}{@{}llrrrr@{}}
\toprule
State & Algorithm & Min centroid-Reock & Mean centroid-Reock & EC & Package \\
\midrule
NC & \metis{} & TBD & TBD & TBD & TBD \\
NC & \cvdgeo{} & TBD & TBD & TBD & TBD \\
NC & \bfsg{} & TBD & TBD & TBD & TBD \\
NC & \mka{} & TBD & TBD & TBD & TBD \\
FL & \metis{} & TBD & TBD & TBD & TBD \\
FL & \mka{} & TBD & TBD & TBD & TBD \\
WA & \metis{} & TBD & TBD & TBD & TBD \\
WA & \mka{} & TBD & TBD & TBD & TBD \\
\bottomrule
\end{tabular}
\end{table}

For publication, the package should record:
\begin{itemize}
\item input graph and centroid hashes;
\item base seed and bisection tree node paths;
\item \texttt{mka\_orientations} and metric;
\item for every split node, the selected angle and score;
\item final plan assignments and validation results;
\item identical scoring scripts for all compared algorithms.
\end{itemize}

\subsection{A Worked Candidate Ledger}

The key reader-facing evidence is a candidate ledger. A small node can be
shown as follows:

\begin{center}
\begin{tabular}{rrrrrr}
\toprule
$\theta$ & prefix pop & left shape & right shape & min score & decision \\
\midrule
$0^\circ$ & 51.2\% & long & compact & 0.31 & reject \\
$45^\circ$ & 50.6\% & blocky & blocky & 0.44 & keep \\
$90^\circ$ & 52.0\% & compact & long & 0.29 & reject \\
$135^\circ$ & 49.8\% & split tail & blocky & 0.37 & reject \\
\bottomrule
\end{tabular}
\end{center}

This is the Moving-Knife equivalent of ``show the split, not say the split.''
The reader can see why one direction wins: it gives both halves acceptable
enclosing-circle geometry.

\subsection{Orientation Granularity}

The default 180-orientation sweep tests one-degree increments. Smaller sweeps
such as 36 orientations can be useful for development, but production claims
should report the granularity explicitly. The benchmark question is not merely
runtime. It is whether coarser scans select the same bisection-tree angles and
produce comparable final validation metrics.

\subsection{Partisan and Legal Metrics}

\mka{} does not read election returns in the split routine. That makes the
structure step formally independent of partisan outcome data, but it does not
prove partisan neutrality of final plans. Recursive bisection, geography, and
population distribution can still create predictable electoral effects.
Publication claims should therefore compare final plans across seeds using the
same partisan metrics used elsewhere in BISECT.
